% Generated draft: requires mathematical and source review.
\documentclass[11pt]{article}
\usepackage[margin=25mm]{geometry}
\usepackage{fontspec}
\setmainfont{Latin Modern Roman}
\usepackage{amsmath,amssymb,mathtools,unicode-math}
\setmathfont{Latin Modern Math}
\usepackage{enumitem,needspace,xurl,hyperref}
\setlength{\parindent}{0pt}
\setlength{\parskip}{5pt}
\setlength{\emergencystretch}{3em}
\newcommand{\N}{\mathbb{N}}
\newcommand{\Z}{\mathbb{Z}}
\newcommand{\Q}{\mathbb{Q}}
\newcommand{\R}{\mathbb{R}}
\newcommand{\C}{\mathbb{C}}
\newcommand{\F}{\mathbb{F}}
\newcommand{\A}{\mathbb{A}}
\newcommand{\PP}{\mathbb{P}}
\newcommand{\E}{\mathbb{E}}
\newcommand{\eps}{\varepsilon}
\newcommand{\dd}{\,\mathrm{d}}
\begin{document}
\section*{UnsolvedMath ID: 20002311}
\textbf{Source problem code:} catalogue entry
\subsection{Definitions and mathematical statement}
Question 9 asks whether, for a quasiminimal excellent class \(\mathcal K\) and parameters \(\lambda<\kappa\),
\[
 \mathcal K\text{ is }\kappa^+\text{-computably categorical}
 \quad\Longrightarrow\quad
 \mathcal K\text{ is }\lambda^+\text{-computably categorical}.
\]
The source does not reproduce the quasiminimal-excellence axioms, cardinal/recursion conventions, presentation domains, permitted parameters, or definition of successor-cardinal computable categoricity. These are not interchangeable with ordinary countable Turing computability. The aid uses the AIM admissible-recursion convention, whose full specification is also absent from this record.

Its special case fixes an integer \(n\ge2\). Let \(\mathcal K_n\) consist of infinite structures in the language of one equivalence relation \(E\), each equivalence class having exactly \(n\) elements, with closure
\[
 \operatorname{cl}(A)=\bigcup\{[a]_E:a\in A\}.
\]
The aid identifies this finite-fiber pregeometry class as quasiminimal excellent. For every successor cardinal \(\theta=\mu^+\) in the stated recursion setting, it reports uniform relative \(\theta\)-computable categoricity: from the diagrams of any two appropriate presentations of the same model one obtains an isomorphism by a uniform \(\Delta_1\) procedure relative to those diagrams. For \(\theta\)-computable copies, oracle elimination gives a \(\theta\)-computable isomorphism without the relative oracles. Thus both relative and nonrelative downward transfers hold within \(\mathcal K_n\); no general transfer for all quasiminimal excellent classes is concluded.

The reported map uses the presentations' prescribed well orders. Choose the least representative of each finite equivalence class, pair the increasing enumerations of representatives, then pair the \(j\)-th elements inside corresponding classes for \(j<n\). On a common infinite initial-ordinal domain, each representative set has the same order type as that domain, so this pairing is well-defined. In the aid's admissible setting these representative and bounded finite-coordinate operations give the claimed \(\Delta_1\) map. The exact relation between a presentation domain and the recursion index \(\theta\) must still be checked against the omitted convention; this description does not silently identify them.

The known fixed class size is part of this construction. The result is a full successor-index spectrum for this elementary finite-fiber example, not a discovery algorithm for arbitrary pregeometries or a solution of the original unrestricted downward-categoricity question.
\subsection{Short English statement}
The question asks whether effective categoricity transfers downward between successor-cardinal recursion levels in quasiminimal excellent classes. The aid proves a finite-equivalence-class example using least representatives and finite fiber coordinates. The general recursion conventions and unrestricted transfer remain outside this supplied special case.
\subsection{Sources}
\begin{itemize}
\item Local UnsolvedMath registry, permanent ID 20002311, exact source code catalogue entry; canonical \texttt{aim-logic-notes.json}, \texttt{notes[86]}, tag \texttt{question}; AIM Workshop Problem Lists, Logic, \emph{Computable stability theory}, Question 9. Created and updated 14 August 2026.
\item Source URL \url{https://aimath.org/pastworkshops/computestabproblems.pdf}. The canonical excerpt contains the downward implication but omits its earlier technical conventions. The PDF was not accessed.
\item Attempt 1: \texttt{partial\_result}, problem status at run \texttt{unknown}, rigor label \texttt{proved}, novelty confidence \texttt{low}; classification \texttt{PARTIAL-PROGRESS}, registry status \texttt{partially\_solved}. Candidate: fixed finite-fiber uniform relative successor-categoricity spectrum. No separate authoritative theorem citation or full admissible-recursion definition is supplied.
\item Ulam AI / UnsolvedMath Contributors; CC BY 4.0. Machine-generated research aid, not peer reviewed; independently verify before citation.
\end{itemize}
\end{document}
